% SPDX-License-Identifier: Apache-2.0
\section{A Line Ahead of the Beam}
\label{sec:x-scanpair}

\S\ref{sec:x-scanout} fills one line from memory and reads it back.
A scanout has to do that for every line of every frame while the beam
moves, so each line must be in place before the beam reaches it, and
when one is not, somebody has to know. Two parts in
\code{scanout} in \code{//lib/parts} do this (issue 151):

\begin{center}\footnotesize
\begin{tabular}{@{}l@{}}
\code{Rom -- AxiPer -- AxiHost -- LineFetch ]--- ChanCdc ---[ LinePair} \\
\code{\phantom{Rom -- AxiPer -- AxiHost -- }ScanFetch ]<-- ChanCdc <--[ LinePair} \\
\hspace{1.1cm}\textit{bus clock}
\hspace{3.6cm}\textit{pixel clock} \\
\end{tabular}
\end{center}

The words go down the first row and the requests come back up the
second.

\code{LinePair}, on the pixel clock, holds two lines. The beam reads
one while the other fills, and they change places at the start of
each line. At that start it asks for the line after the one about to
be shown, providing the full duration of the previous line for arrival.
The address is a running sum, a stride a line, so nothing multiplies.
The frame's base is a register taken at the vertical sync, so a host
that draws in one buffer and shows another flips them with one write.
A column the beam reads before its word arrived sets \code{under}, a
sticky bit a host reads and clears, which is how a run on the board
shows that no line ever starved. \code{ScanFetch}, on the bus clock,
takes each request across and starts \code{LineFetch} on it.

\lstinputlisting[rangebeginprefix=//\ begin\{,rangebeginsuffix=\},%
rangeendprefix=//\ end\{,rangeendsuffix=\},includerangemarker=false,%
linerange=pair_run-pair_run,caption={\code{lib/parts/src/scanout.rs}:
the pixel side.},label={lst:x-scanpair-pair}]{lib/parts/src/scanout.rs}

\subsection{What a line must survive}

On the board the memory is AMD's MIG in \code{//ddr3}, for an
MT41K256M16 at 2.5\,ns a clock. Its refresh takes $t_{RFC}$ =
260\,ns once every $t_{REFI}$ = 7.8\,$\mu$s, and a bank
miss costs $t_{RP} + t_{RCD}$ = 27.5\,ns before a CAS latency of
six clocks. The video mode is 640 by 480, 31.8\,$\mu$s a line,
and a visible line is 640 words, forty bursts of sixteen. Since a
line is fetched during the one before it, what it must survive is
forty bursts and the four refreshes that land in a line, against the
line's time, not one burst's latency against one pixel.

\subsection{What the run checks}

The raster is the testbench's, eight visible columns of sixty-four
and four visible rows of six, on a pixel clock unrelated to the bus
clock. The memory takes a \code{hold} that stalls its beats, standing
in for a refresh. Two frames, with a short stall in every line, show
every pixel as the word at its place in the frame, and the underflow
bit stays clear. A later frame stalls the memory for an entire
line, so the next line cannot arrive in time: the bit sets and stays
set until a write clears it. Both units are checked against the run
under nvc and Verilator.

\begin{figure}[htbp]
\centering
\input{scanpair_timing.tex}
\caption{The long stall: \code{start} does not come back while the
memory holds the row after, \code{rgot} is short at the next line's
start, and \code{starved} sets and stays set.}
\label{fig:x-scanpair}
\end{figure}

\lstinputlisting[caption={\code{ex\_scanpair.rs}. Run with \code{bazel run //lib/examples:ex\_scanpair}.},label={lst:x-scanpair}]{lib/examples/ex_scanpair.rs}

\lstinputlisting[style=ascii,caption={What \code{ex\_scanpair} printed when the build ran it: the pixels checked, the underflow bit at each frame's end, and the netlist of \code{LinePair}.},label={lst:x-scanpair-out}]{out_scanpair.txt}
